\section{Why Response Determines Utility}
\label{sec:theory}

Throughout, $f_A(x)$ is the prediction of a predictor that consumes the anchor
context and the selected set $A$, $y$ is the outcome, $\loss$ is the task loss,
and $\dAj(x)=f_{A\cup\{j\}}(x)-f_A(x)$ is the predictor's response to candidate
$j$ at state $A$. The marginal utility of that candidate is
\begin{equation}
 \marginal(j\mid A)=\E\!\left[\loss(f_A,y)-\loss(f_{A\cup\{j\}},y)\right]-\lambda c_j .
 \label{eq:marginal-utility}
\end{equation}
The predictor is frozen while utilities are generated, so every statement below
is conditional on the state and on the predictor.

\subsection{A two-sided bound for smooth losses}

\begin{proposition}[Response expansion]
\label{prop:expansion}
Assume that for every episode $f\mapsto\loss(f,y)$ is twice continuously
differentiable on the segment between $f_A(x)$ and $f_{A\cup\{j\}}(x)$, and that
$\|\grad^2\loss(z,y)\|_2\leq L$ on that segment. Then
\begin{align}
 \marginal(j\mid A)
 &=
 -\E\!\left[\grad\loss(f_A,y)^\top \dAj\right]
 -\tfrac12\E\!\left[\dAj^\top\grad^2\loss(\xi,y)\dAj\right]
 -\lambda c_j,
 \label{eq:prop-expansion}\\
 \left|\marginal(j\mid A)
 +\E\!\left[\grad\loss(f_A,y)^\top \dAj\right]+\lambda c_j\right|
 &\leq
 \frac{L}{2}\E\|\dAj\|^2 .
 \label{eq:prop-bound}
\end{align}
\end{proposition}

The leading term is the alignment between the loss gradient at the current
prediction and the direction in which the candidate moves the predictor. The
remainder is second order in the response magnitude: the first-order term is
accurate exactly when the candidate perturbs the predictor little, and a large
response can cost at most $(L/2)\E\|\dAj\|^2$. The response is observable at
decision time; only the gradient is not.

\subsection{An exact identity for Bregman losses}

\begin{proposition}[Exact identity]
\label{prop:bregman}
Let $\loss(f,y)=D_\phi(y,f)=\phi(y)-\phi(f)-\grad\phi(f)^\top(y-f)$ for convex
differentiable $\phi$ --- squared loss, deviance losses, and exponential-family
negative log-likelihoods in their natural parametrisation. Then
\begin{equation}
 \marginal(j\mid A)
 =
 \E\!\left[\left(\grad\phi(f_{A\cup\{j\}})-\grad\phi(f_A)\right)^\top(y-f_A)\right]
 -\E D_\phi\!\left(f_A,f_{A\cup\{j\}}\right)
 -\lambda c_j .
 \label{eq:prop-bregman}
\end{equation}
\end{proposition}

Equation~\eqref{eq:prop-bregman} is exact and requires no smoothness constant.
The second term is a non-positive movement penalty in the Bregman geometry; the
first is the inner product between the change in the loss gradient and the
current residual. With $\phi(z)=\|z\|^2$ it becomes
$\E[2\dAj(y-f_A)-\|\dAj\|^2]-\lambda c_j$, the squared-loss expression used by
response-aware routers, which is therefore a special case rather than the
whole story.

\subsection{Explicit constants for cross-entropy}

\begin{proposition}[Softmax cross-entropy]
\label{prop:ce}
For logits $f\in\R^C$, one-hot $y$, and softmax cross-entropy
$\loss(f,y)=-f_y+\log\sum_k e^{f_k}$ with $p=\softmax(f)$, the Hessian
$\grad^2\loss=\diag(p)-pp^\top$ is positive semidefinite with spectral norm at
most $1/2$. Consequently
\begin{equation}
 \langle e_y-p_A,\dAj\rangle-\tfrac14\E\|\dAj\|^2-\lambda c_j
 \;\leq\;
 \marginal(j\mid A)
 \;\leq\;
 \langle e_y-p_A,\dAj\rangle-\lambda c_j .
 \label{eq:prop-ce}
\end{equation}
\end{proposition}

The upper bound holds because the second-order term penalises any positive
semidefinite curvature; the lower bound says that a large response can lose at
most a quarter of its squared size in logit space. Both constants are explicit,
which matters when the response is used as a feature.

\subsection{Response sufficiency}

\begin{proposition}[Response sufficiency]
\label{prop:suff}
For squared loss, conditioning on the state gives
\begin{equation}
 \E\!\left[\widetilde{\marginal}(j\mid A)\mid x,A\right]
 =
 \langle g_A(x),\dAj(x)\rangle-\|\dAj(x)\|^2-\lambda c_j,
 \qquad
 g_A(x)=2(\mu(x)-f_A(x)),
 \label{eq:prop-suff}
\end{equation}
where $\mu(x)=\E[y\mid x,A]$.
\end{proposition}

Conditional on the state, the Bayes-optimal squared-loss utility is affine in
the response vector plus a known quadratic penalty. State features supply the
coefficient $g_A(x)$, response features supply $\dAj$, and the penalty is a
function of the response alone. Any score that is independent of the response
can only recover the state average and provably omits the interaction
$\langle g_A,\dAj\rangle$ that state-conditioned selection exists to exploit.

\subsection{The identifiability corollary}

\begin{corollary}[What is left to learn]
\label{cor:identifiability}
Under \cref{prop:suff}, the only unknown in the optimal score is
$g_A(x)=2(\mu(x)-f_A(x))$, the residual of the fixed predictor. In particular,
if the predictor is calibrated so that $\E[y-f_A(x)\mid x,A]=0$ for every
observable state, the conditional expected utility of every candidate collapses
to $-\|\dAj(x)\|^2-\lambda c_j$, which is non-positive and decreasing in the
response magnitude.
\end{corollary}

The corollary is the paper's central practical statement. A better predictor is
one whose residual is closer to unpredictable, so the term the router must
estimate shrinks towards zero while the curvature penalty remains. The oracle
opportunity --- what a router could achieve if it knew the outcomes --- does not
shrink; what shrinks is the part that is identifiable before the outcome is
observed. This is why strong predictors are not simply a harder version of the
same problem.

\subsection{From estimation error to routing regret}

\begin{theorem}[One-step routing regret]
\label{thm:onestep}
At a fixed state $A$, suppose the utility estimator satisfies
$|\widehat{\marginal}(j\mid A)-\marginal(j\mid A)|\leq\epsilon$ for every
remaining candidate. If $\widehat{j}$ maximises the estimate and $j^*$ maximises
the true utility, then
$\marginal(j^*\mid A)-\marginal(\widehat{j}\mid A)\leq 2\epsilon$.
\end{theorem}

Combining \cref{thm:onestep} with \cref{prop:expansion} bounds the regret of a
response-based router by twice the estimation error of its first-order surrogate
plus the second-order response term. The bound is per decision; sequential
composition additionally requires that locally accurate choices do not visit a
qualitatively different state, which is why \cref{sec:diagnostics} measures
identifiability directly instead of assuming it.
